\section*{Resumen}
\addcontentsline{toc}{section}{Resumen}
La minería ilegal en la Amazonía colombiana genera afectaciones ambientales y sociales que requieren mecanismos automatizados de vigilancia. Los videos capturados por las cámaras FLIR aerotransportadas de la Fuerza Aeroespacial Colombiana permiten observar con detalle zonas de difícil acceso, pero su análisis automático enfrenta tres dificultades: la simbología de vuelo (HUD) superpuesta y el ruido de la imagen, la alta correlación entre fotogramas de un mismo video y la limitación de las cajas delimitadoras para representar regiones extensas como ríos y zonas de alteración del suelo. Este proyecto implementa un pipeline funcional y reproducible que articula tres componentes: (1) la limpieza de los videos, que elimina el HUD mediante segmentación heurística e inpainting de video y reduce el ruido con métodos autosupervisados escogidos según el régimen de operación del sensor; (2) un particionamiento basado en agrupamiento de fotogramas altamente correlacionados, que evita la fuga de datos entre entrenamiento, validación y prueba; y (3) una arquitectura de segmentación panóptica (Panoptic FCN) con supervisión puntual, que representa conjuntamente regiones ambientales y objetos asociados con actividad minera. En la etapa de limpieza, el HUD se redujo en una mediana de 98,4\,\% a 98,6\,\% de sus píxeles por imagen en los tres videos, y la caracterización del ruido mostró que este depende del régimen del sensor y no solo del video. \pendiente{resultados finales de reducción de ruido, particionamiento y segmentación panóptica.}
